\documentclass[UTF8,a4paper,11pt]{ctexart}
\usepackage[left=2.6cm,right=2.6cm,top=2.8cm,bottom=2.8cm]{geometry}
\usepackage{amsmath,amssymb,amsthm,bm,mathtools}
\usepackage{booktabs}
\usepackage{array}
\usepackage{multirow}
\usepackage{enumitem}
\usepackage{hyperref}
\usepackage{graphicx}
\usepackage{setspace}
\usepackage{indentfirst}
\usepackage{xcolor}
\usepackage{longtable}

\setlength{\parindent}{2em}
\linespread{1.2}
\hypersetup{hidelinks}

\newtheorem{theorem}{定理}[section]
\newtheorem{proposition}{命题}[section]
\newtheorem{remark}{注}[section]
\newtheorem{definition}{定义}[section]
\newtheorem{assumption}{假设}[section]

\title{面向 Chrono NSC 的多 Patch / 多 Subpatch 保持 Wrench 等效的 SDF 接触约化\\中文论文骨架（方法导向草稿）}
\author{}
\date{}

\begin{document}
\maketitle

\begin{abstract}
针对 Chrono 的非光滑接触（non-smooth contact, NSC）框架在高精度连续区域接触建模中的局限，本文提出一种面向 SDF 几何表征的多 patch / 多 subpatch 保持 wrench 等效的接触约化方法。该方法以符号距离场为几何基础，首先构造高密度接触区域参考模型；随后依据连通性、法向变化、几何曲率与 wrench 误差将接触区域自适应分解为多个 patch 与 subpatch；最后针对每个 subpatch 构建少量代表点组成的 reduced support set，使其在给定误差阈值下逼近 dense 接触模型的可行 wrench 锥，并将这些代表点重新注入 Chrono NSC 的标准点接触求解链路中。与传统基于少量缓存点的接触流形维护方式不同，本文方法不是简单进行几何删点，而是将接触约化重新表述为“由 dense admissible wrench cone 到 reduced admissible wrench cone 的误差控制投影”。在理论上，本文给出了局部共面 subpatch 的三点 wrench 完备性、法向可行性与中心压力点支撑域之间的几何条件，并提出一套以 wrench 误差、CoP 误差、内部 gap 误差和拓扑误差共同驱动的自适应细化策略。本文进一步给出面向凸轮连续区域接触、点接触以及长条形接触（齿轮、插销--轴承等）的统一实验设计框架。该方法的目标并非逐点重建真实压力场，而是在保持 Chrono NSC 全局互补求解框架不变的前提下，提高连续接触区域在刚体动力学层面的合力、合力矩与作用线等效精度，为高保真 SDF 接触与现有多体动力学求解器的耦合提供一条可落地的技术路线。
\end{abstract}

\noindent\textbf{关键词：}符号距离场；非光滑接触；Chrono；接触流形；wrench 等效；多体动力学

\tableofcontents
\clearpage

\section{引言}

刚体接触动力学是多体系统仿真中的核心问题之一，也是机器人、机械系统、图形学仿真与数字样机等领域的共同基础。对于包含频繁接触与摩擦的系统，工业界和学术界广泛采用基于互补约束、速度层非穿透条件以及库仑摩擦的非光滑接触框架。Chrono 的 NSC 体系是其中具有代表性的实现之一，其优势在于建模清晰、求解流程成熟、适合处理复杂刚体系统中的大量接触与碰撞。然而，当接触由局部点接触扩展到连续区域接触、长条形接触、多 patch 接触或者由几何高阶特征主导的复杂接触时，传统点接触驱动的流形表示很难同时兼顾精度、稳定性与效率。

这一问题的根源并不只是“接触点数量不够”，而在于连续接触区域在力学上携带的信息远比少量几何接触点更加丰富。一个真实接触区域包含接触拓扑、法向变化、接触面积尺度、局部滑移分布以及接触作用线等多方面信息。若直接将其压缩为少数离散点并送入互补求解器，则虽然仍可能维持非穿透约束的可解性，但其 body-level 的合力、合力矩乃至作用线位置都可能产生显著偏差。尤其在凸轮接触、面--面压靠、齿轮啮合、轴承--插销和线接触扩展为长条接触等情形中，这种误差往往会累积为明显的刚体响应误差。

另一方面，符号距离场（signed distance field, SDF）为接触建模提供了新的几何基础。相较于仅以表面三角片或离散采样点表示几何，SDF 能够直接提供局部 gap、法向和几何邻域信息，因此更适合构造连续接触区域及其微分结构。在此基础上，如果仍沿用传统“接触点缓存 + 点接触求解”的思路，则 SDF 的优势难以充分发挥。更合理的路线应当是：先由 SDF 构造高保真的接触区域参考模型，再将其投影为与现有求解器兼容的低维接触基。

已有研究表明，对单个凸接触区域，可通过等效接触点（Equivalent Contact Point, ECP）与接触 wrench 的方式统一描述点接触、线接触和面接触；也有研究采用基于接触曲面和压力分布积分的方法，直接由连续接触区域求得结果力与结果力矩。受这些思想启发，本文不再把接触约化视为一个“如何挑几个代表点”的几何问题，而是将其重新表述为一个“如何保持连续接触区域在 wrench 空间中的可行集”的力学投影问题。

基于上述认识，本文提出一种面向 Chrono NSC 的多 patch / 多 subpatch 保持 wrench 等效的 SDF 接触约化方法。该方法强调以下三点：第一，以 SDF 为基础构造 dense 接触区域参考模型；第二，以局部可行 wrench 锥而非几何点集作为约化对象；第三，以误差驱动的多 patch / 多 subpatch 细化替代固定规则的流形删点。通过这一设计，本文希望在不重写 Chrono NSC 全局互补求解框架的前提下，显著提高连续接触区域在刚体动力学层面的等效精度。

本文的主要贡献可概括如下：
\begin{enumerate}[label=(\arabic*)]
    \item 提出一种面向 SDF 接触区域的多 patch / 多 subpatch wrench-preserving reduction 框架，将接触约化从几何删点提升为对可行 wrench 锥的误差控制投影；
    \item 给出局部共面 subpatch 的三点 wrench 完备性分析，并建立法向可行性与中心压力点（CoP）支撑域之间的几何判据；
    \item 构建一套融合几何误差、wrench 锥误差、CoP 误差、内部 gap 误差与拓扑误差的自适应 refinement 机制；
    \item 给出与 Chrono NSC 集成的实现路径，使所提方法能够在保持现有全局互补求解器不变的条件下落地。
\end{enumerate}

本文余下内容安排如下：第~\ref{sec:related} 节回顾相关工作；第~\ref{sec:preliminaries} 节给出问题背景与基本定义；第~\ref{sec:method} 节详细介绍本文方法；第~\ref{sec:analysis} 节给出理论分析与误差讨论；第~\ref{sec:experiments} 节给出实验设计；第~\ref{sec:conclusion} 节总结全文并讨论未来工作。

\section{相关工作}
\label{sec:related}

\subsection{基于互补约束的非光滑接触动力学}

非光滑接触动力学通常通过非穿透约束、速度层互补条件和库仑摩擦律建立模型。在多体系统中，这类方法具有较好的鲁棒性与工程可落地性，特别适合刚体碰撞、频繁接触和较大时间步长的仿真。Chrono 的 NSC 路线即属于该范式，其核心思想是将接触约束装配到系统描述器中，通过互补求解器或相关变体求解接触反力与系统速度更新。该类方法的主要优点是求解框架统一、适合复杂刚体系统耦合，但其接触层通常依赖离散接触点表示，因此对于连续区域接触的高精度建模存在天然张力。

\subsection{基于接触点缓存的流形维护}

在工程实现上，常见碰撞后端会在窄相检测之后维护少量持久接触点组成的接触流形，以提高时间连续性、暖启动效果和求解效率。该类方法的接触流形本质上更接近“接触点缓存”而非“连续接触区域的物理近似”。它们通常以启发式方式保留最深点、覆盖面积较大的点或时间上持续存在的点，因而在点接触和局部小面积接触中效果较好，但在长条接触、复杂多 patch 接触和大尺度连续区域接触中，可能出现接触作用线漂移、力矩失真与拓扑混叠。

\subsection{基于接触曲面和压力场的分布式接触}

与离散点流形不同，另一类方法从接触曲面出发，在连续接触区域上定义法向压力、切向牵引或相关势函数，再通过面积分得到结果力和结果力矩。此类方法更强调连续区域本身的几何与力学意义，因此在需要高精度 body-level wrench 或接触面演化分析时更具优势。然而，其代价往往是求解模型更复杂、与现有刚体互补求解框架的直接兼容性较弱。

\subsection{等效接触点与接触 wrench 表示}

等效接触点（ECP）相关工作指出，对单个凸接触区域，可以用一个等效接触点及其对应的接触 wrench 统一表征点接触、线接触与面接触。这一思想的重要启发在于：连续区域接触并不必然要求大量离散点才能在刚体动力学层面被精确表征；关键在于是否保留了正确的结果 wrench 与作用线信息。然而，直接使用 patch-level 六维接触元又需要改写现有求解器的数据结构与约束装配过程。

\subsection{SDF 几何与接触建模}

SDF 在接触建模中的优势主要体现在三个方面：其一，可直接提供局部 gap 与法向；其二，便于构造连续接触区域及其局部微分结构；其三，适合与体素、稀疏体数据结构和自适应采样耦合。现有基于 SDF 的接触研究多集中于穿深估计、法向计算、柔顺势函数构造和图形学仿真。相比之下，如何在 SDF 基础上面向刚体 NSC 框架构建“高保真但仍可低维求解”的接触约化，仍缺乏系统研究。

\subsection{本文方法与现有工作的区别}

本文方法既不同于传统少量缓存点的流形维护，也不同于直接在全局系统中引入 patch-level 六维接触元。本文的出发点是：以 SDF dense 接触区域为参考对象，在局部 patch / subpatch 层面保持其可行 wrench 锥和当前参考 wrench，并最终将这一等效结果重新投影回标准点接触集合。换言之，本文不是在求解器末端进行“删点”，而是在接触建模接口层进行“保持 wrench 等效的低维投影”。这使其既能显著改善连续区域接触的刚体层面等效精度，又能保持对 Chrono NSC 现有求解链路的兼容性。

\section{预备知识与问题背景}
\label{sec:preliminaries}

\subsection{Chrono NSC 中的接触求解链路}

在 Chrono NSC 框架中，碰撞后端首先给出候选接触对及其对应的接触点信息；随后这些信息被封装为接触容器中的接触对象，并进一步装配到全局系统描述器中；最终，互补求解器在变量增量与约束乘子空间中求得各接触反力分量。该链路天然更偏向“一个离散接触点对应一个接触约束块”的建模模式，因此若直接输入少量离散接触点，则求解器能高效工作，但连续区域接触的物理语义会在进入求解器之前被压缩。

\subsection{为何仅增加接触点数量并不能根本解决问题}

直观上，似乎可以通过简单增加接触点数量来提高精度。然而，这种做法存在两方面局限：第一，离散点的几何分布未必能正确保持结果力矩与作用线；第二，接触点数量的增加会直接增大全局求解规模，却不一定显著改善等效 wrench 精度。因此，问题的关键不是“输入多少点”，而是“输入的点是否构成了对连续区域接触 wrench 的良好低维基”。

\subsection{基本记号}

为方便后续推导，本文统一采用如下记号：
\begin{itemize}[leftmargin=2em]
    \item $\Omega$：某一候选接触对的连续接触区域；
    \item $\Omega_\ell$：第 $\ell$ 个 patch；
    \item $\Omega_{\ell r}$：第 $\ell$ 个 patch 内的第 $r$ 个 subpatch；
    \item $\mathcal{Q}$：dense 积分点集合；
    \item $Q_{\ell r}$：subpatch 的 reduced support set；
    \item $w$：局部六维 wrench，按 $(T_x,T_y,N,M_x,M_y,M_z)^\top$ 排列；
    \item $\mathcal{C}_{\mathrm{dense}}$ 与 $\mathcal{C}_{\mathrm{red}}$：分别表示 dense 与 reduced 可行 wrench 锥；
    \item CoP：中心压力点（center of pressure）；
    \item $e_{\mathrm{plane}}, e_{\mathcal{C}}, e_w, e_{\mathrm{CoP}}, e_g$：几何与力学误差指标。
\end{itemize}

\section{方法：面向 Chrono NSC 的多 Patch / 多 Subpatch 保持 Wrench 等效的 SDF 接触约化}
\label{sec:method}

\subsection{方法概述}

现有刚体接触仿真中，基于少量接触点缓存的离散接触流形更强调时间连续性与求解效率，而非连续接触区域的高保真力学表征。与之相对，基于接触曲面与压力场积分的思路体现了“surface-first”的建模思想：先构造接触面，再由接触面上的分布量得到合力与合力矩。进一步地，等效接触点（Equivalent Contact Point, ECP）类方法表明，对单个凸接触区域，可以用一个等效接触点及其对应的接触 wrench 来统一线接触与面接触。上述认识共同说明：高精度刚体接触的关键，不在于缓存少量几何点，而在于如何在刚体动力学层面保留连续接触区域的等效 wrench。

基于此，本文提出一种面向 Chrono NSC 的\textbf{多 Patch / 多 Subpatch 保持 Wrench 等效的 SDF 接触约化方法}。该方法以 SDF 为几何基础，首先为每个候选接触对构造高密度接触区域参考模型；随后将连续接触区域分解为多个 patch 与 subpatch；最后对每个 subpatch 构造少量代表点组成的 reduced support set，使其在给定误差阈值下逼近 dense 接触模型所张成的可行 wrench 锥，并将所得代表接触点重新注入 Chrono NSC 接触容器。该方法定位于碰撞--接触建模接口层，而不修改全局 NSC 互补求解框架。

本文方法的核心思想可概括为：
\begin{enumerate}[label=(\arabic*)]
    \item \textbf{surface-first}：先由 SDF 构造高密度接触区域，而非直接信任离散点流形；
    \item \textbf{wrench-first}：约化目标不是保留若干几何点，而是逼近 dense 接触模型的可行 wrench 锥；
    \item \textbf{adaptive refinement}：以几何误差、wrench 误差、CoP 误差与连续区域 gap 误差共同驱动自适应细分；
    \item \textbf{solver-compatible}：最终输出仍为普通点接触，从而直接复用 Chrono NSC 的求解链路。
\end{enumerate}

\subsection{问题定义}

考虑两个刚体 $A$ 与 $B$，其几何分别由符号距离场 $\phi_A(x)$ 与 $\phi_B(x)$ 描述。对于任一候选接触对，在相对位姿给定时构造局部 gap 场 $g(x)$。在本文中，$g(x)$ 可由双 SDF 组合、单边表面采样或其他等价构造获得，其目标是给出局部非穿透约束的几何刻画。接触区域定义为
\begin{equation}
\Omega = \{x \in \mathbb{R}^3 \mid g(x) \le 0\}.
\end{equation}

为获得高精度参考接触模型，将 $\Omega$ 离散为带面积权重的高密积分点集合
\begin{equation}
\mathcal{Q} = \{(x_k, a_k, n_k)\}_{k=1}^{N_q},
\end{equation}
其中 $x_k$ 为接触积分点，$a_k$ 为面积权重，$n_k$ 为由 SDF 梯度给出的局部单位法向，即
\begin{equation}
 n_k = \frac{\nabla \phi(x_k)}{\|\nabla \phi(x_k)\|}.
\end{equation}

在每个积分点上，引入局部接触力变量
\begin{equation}
\lambda_k =
\begin{bmatrix}
 t_{k1} \\
 t_{k2} \\
 n_k^c
\end{bmatrix},
\qquad
n_k^c \ge 0,
\qquad
\sqrt{t_{k1}^2 + t_{k2}^2} \le \mu_k n_k^c,
\end{equation}
其中 $n_k^c$ 表示法向反力标量，$(t_{k1}, t_{k2})$ 表示局部切向反力分量，$\mu_k$ 为局部摩擦系数。

对于参考点 $c$，记 $x_k-c$ 在局部切平面基 $(e_1,e_2,e_3)$ 中的坐标为 $(u_k,v_k,0)$。则每个积分点对局部六维 wrench 的映射矩阵可写为
\begin{equation}
G(x_k) =
\begin{bmatrix}
1 & 0 & 0 \\
0 & 1 & 0 \\
0 & 0 & 1 \\
0 & 0 & v_k \\
0 & 0 & -u_k \\
-v_k & u_k & 0
\end{bmatrix}.
\end{equation}

因此，dense 接触模型可产生的任意局部 wrench 为
\begin{equation}
 w = \sum_{k=1}^{N_q} G(x_k)\lambda_k.
\end{equation}

进一步定义 dense 可行 wrench 锥
\begin{equation}
\mathcal{C}_{\mathrm{dense}}
=
\left\{
\sum_{k=1}^{N_q} G(x_k)\lambda_k
\;\middle|\;
 n_k^c \ge 0,\;
 \|(t_{k1},t_{k2})\|_2 \le \mu_k n_k^c
\right\}.
\end{equation}

本文约化的目标不是保留 $\mathcal{Q}$ 的全部几何点，而是构造一个低维 reduced support set，使其对应的 reduced wrench 锥尽可能逼近 $\mathcal{C}_{\mathrm{dense}}$。

\subsection{多 Patch / 多 Subpatch 分解}

\subsubsection{Patch 分解}

由于一个接触对可能包含多个不连通接触区域，本文首先在高密积分点集合 $\mathcal{Q}$ 上建立邻接图
\begin{equation}
(i,j)\in E
\iff
\|x_i - x_j\| < h
\;\land\;
n_i \cdot n_j > \cos\theta_0
\;\land\;
g(x_i)\le 0,\; g(x_j)\le 0,
\end{equation}
其中 $h$ 为局部采样尺度，$\theta_0$ 为法向相容阈值。图的连通分量定义为接触 patch：
\begin{equation}
\Omega = \bigcup_{\ell=1}^{N_p} \Omega_\ell.
\end{equation}

这种处理确保不同的连通接触区域不会被错误合并为单一 patch。

\subsubsection{Subpatch 分解}

仅按连通性分块仍不足以保证局部 wrench 等效的精度。为此，本文进一步将每个 patch $\Omega_\ell$ 分解为若干 subpatch：
\begin{equation}
\Omega_\ell = \bigcup_{r=1}^{m_\ell} \Omega_{\ell r}.
\end{equation}

subpatch 的细分由如下误差指标联合驱动：
\begin{enumerate}[label=(\roman*)]
    \item \textbf{局部平面化误差}
    \begin{equation}
    e_{\mathrm{plane}} = \kappa_{\max} D,
    \end{equation}
    其中 $D$ 为当前块的直径，$\kappa_{\max}$ 为最大主曲率；

    \item \textbf{二阶矩误差}
    \begin{equation}
    e_{\Sigma}
    =
    \frac{\|\Sigma_{\mathrm{red}}-\Sigma_{\mathrm{dense}}\|_F}
    {\|\Sigma_{\mathrm{dense}}\|_F+\varepsilon};
    \end{equation}

    \item \textbf{可行 wrench 锥误差}
    \begin{equation}
    e_{\mathcal{C}}
    =
    \max_{s\in\mathcal{S}_w}
    \frac{|h_{\mathrm{dense}}(s)-h_{\mathrm{red}}(s)|}
    {h_{\mathrm{dense}}(s)+\varepsilon},
    \end{equation}
    其中 $h_{\mathcal{C}}(s)=\sup_{w\in\mathcal{C}} s^\top w$ 为 support function；

    \item \textbf{内部漏穿透误差}
    \begin{equation}
    e_g = \max_{s_i\in \mathcal{S}_{\ell r}} \max(0,-g(s_i)),
    \end{equation}
    其中 $\mathcal{S}_{\ell r}$ 为布置于 subpatch 内部的 sentinel 点集合。
\end{enumerate}

一旦任一误差超过预设阈值，当前块将继续沿主方向、法向变化方向或局部骨架方向进行细分。由此，本文形成由误差控制的 patch-subpatch 自适应层级结构。

\subsection{局部保持 Wrench 等效的 Reduced Support Set}

\subsubsection{局部坐标系与 Reduced Wrench 锥}

对于任意 subpatch $\Omega_{\ell r}$，定义其局部参考点 $c_{\ell r}$ 为面积加权质心，局部法向 $e_3$ 为面积加权平均法向，$e_1,e_2$ 为张成局部切平面的正交基。设 reduced support set 为
\begin{equation}
Q_{\ell r}=\{q_j=(u_j,v_j,0)\}_{j=1}^{m},
\end{equation}
其中 $m$ 为当前 subpatch 的代表点数。对应的 reduced 点力变量为
\begin{equation}
\hat{\lambda}_j=
\begin{bmatrix}
\hat{t}_{j1}\\
\hat{t}_{j2}\\
\hat{n}_j
\end{bmatrix},
\qquad
\hat{n}_j\ge0,
\qquad
\sqrt{\hat{t}_{j1}^2+\hat{t}_{j2}^2}\le \mu_j \hat{n}_j.
\end{equation}

则 reduced 模型对应的可行 wrench 锥为
\begin{equation}
\mathcal{C}_{\mathrm{red}}(Q_{\ell r})
=
\left\{
\sum_{j=1}^{m} G(q_j)\hat{\lambda}_j
\right\}.
\end{equation}

我们的目标是使 $\mathcal{C}_{\mathrm{red}}$ 在给定容忍度下逼近 $\mathcal{C}_{\mathrm{dense}}$。

\subsubsection{三点局部 Wrench 完备性}

当 $\Omega_{\ell r}$ 局部近似共面时，三个非共线代表点已足以生成一个完整的局部六维 wrench 基。具体地，总 wrench 分量满足
\begin{equation}
\begin{aligned}
T_x &= \sum_{j=1}^{m} \hat{t}_{j1},\\
T_y &= \sum_{j=1}^{m} \hat{t}_{j2},\\
N   &= \sum_{j=1}^{m} \hat{n}_j,\\
M_x &= \sum_{j=1}^{m} v_j \hat{n}_j,\\
M_y &= -\sum_{j=1}^{m} u_j \hat{n}_j,\\
M_z &= \sum_{j=1}^{m}\left(u_j \hat{t}_{j2}-v_j \hat{t}_{j1}\right).
\end{aligned}
\end{equation}

对 $m=3$ 且三点非共线的情形，法向分量满足
\begin{equation}
\begin{bmatrix}
1 & 1 & 1\\
v_1 & v_2 & v_3\\
-u_1 & -u_2 & -u_3
\end{bmatrix}
\begin{bmatrix}
\hat{n}_1\\
\hat{n}_2\\
\hat{n}_3
\end{bmatrix}
=
\begin{bmatrix}
N\\
M_x\\
M_y
\end{bmatrix},
\end{equation}
切向分量满足
\begin{equation}
\begin{bmatrix}
1 & 0 & 1 & 0 & 1 & 0\\
0 & 1 & 0 & 1 & 0 & 1\\
-v_1 & u_1 & -v_2 & u_2 & -v_3 & u_3
\end{bmatrix}
\begin{bmatrix}
\hat{t}_{11}\\
\hat{t}_{12}\\
\hat{t}_{21}\\
\hat{t}_{22}\\
\hat{t}_{31}\\
\hat{t}_{32}
\end{bmatrix}
=
\begin{bmatrix}
T_x\\
T_y\\
M_z
\end{bmatrix}.
\end{equation}

当三点非共线时，上述两个矩阵均满秩，因此在忽略库仑锥约束的情况下，三点已对局部六维 wrench 完备。由此可知，\textbf{三点是局部共面 subpatch 的最小 wrench 完备基}。若摩擦可行性不足或几何条件恶化，则提高 $m$ 或继续细分 subpatch。

\subsubsection{法向可行性与 CoP 支撑域}

对 $N>0$，定义中心压力点（center of pressure, CoP）
\begin{equation}
 u_c = -\frac{M_y}{N},
 \qquad
 v_c = \frac{M_x}{N}.
\end{equation}

令 $\lambda_j = \hat{n}_j/N$，则有
\begin{equation}
\lambda_1+\lambda_2+\lambda_3=1,
\qquad
\sum_{j=1}^{3}\lambda_j q_j = (u_c,v_c,0).
\end{equation}

\begin{proposition}
对于三点 reduced support set，存在 $\hat{n}_j\ge 0$ 使法向合力与法向力矩精确匹配，当且仅当 CoP 落在三角形 $\mathrm{conv}\{q_1,q_2,q_3\}$ 内。
\end{proposition}

该命题给出了一个重要设计原则：代表点三角形并非任意选取，而必须覆盖 dense subpatch 在给定工况下可能出现的 CoP 活动区域。

\subsection{代表点位置优化}

\subsubsection{几何初始化}

首先计算 dense subpatch 投影到局部平面的二维点集 $\xi_k$ 的一、二阶矩：
\begin{equation}
A = \sum_k a_k,
\qquad
\bar{\xi} = \frac{1}{A}\sum_k a_k \xi_k,
\qquad
\Sigma = \frac{1}{A}\sum_k a_k(\xi_k-\bar{\xi})(\xi_k-\bar{\xi})^\top.
\end{equation}

令 $LL^\top=\Sigma$。对初始三点，可取标准二维 simplex 顶点 $s_j$ 构造
\begin{equation}
 q_j^{(0)} = \bar{\xi} + \gamma L s_j,
\end{equation}
其中 $\gamma$ 为尺度因子。该初始化使 reduced support set 在均值与二阶尺度上与 dense patch 对齐。

\subsubsection{精度优先目标函数}

最终代表点通过如下优化求得：
\begin{equation}
\min_{Q\subset \mathcal{H}}
\;
\alpha \|c(Q)-\bar{\xi}\|^2
+\beta \|\Sigma(Q)-\Sigma\|_F^2
+\gamma E_{\mathcal{C}}(Q)^2
-\eta \log\det(B_n B_n^\top)
-\zeta \log\det(B_t B_t^\top),
\end{equation}
其中：
\begin{itemize}[leftmargin=2em]
    \item $\mathcal{H}$ 为当前 subpatch 的几何支撑域；
    \item $c(Q)$ 与 $\Sigma(Q)$ 分别为代表点集合的一、二阶矩；
    \item $E_{\mathcal{C}}(Q)$ 为 dense/reduced wrench 锥的 support-function 误差；
    \item $B_n$ 与 $B_t$ 分别为法向与切向矩阵。
\end{itemize}

其中前两项约束 reduced support set 的几何尺度，第三项直接最小化 wrench 锥误差，后两项提升法向与切向矩阵的条件数，避免代表点塌缩或杠杆臂退化。

\subsection{局部参考 Wrench 预测与 Reduced 力分配}

\subsubsection{局部参考 Wrench 预测}

仅逼近可行 wrench 锥仍不足以决定当前时间步应落在锥中的哪一点。因此，本文为每个 dense subpatch 构造一个\textbf{局部参考 wrench 预测器}，得到当前时刻参考 wrench $w_{\ell r}^\star$。

设当前相对 twist 为 $(v_{\mathrm{rel}}, \omega_{\mathrm{rel}})$。本文采用局部 dense 微求解方式，在 dense 积分点上解一个小型 SOCP/QP：
\begin{equation}
\min_{\lambda}
\frac{1}{2}\lambda^\top H \lambda + g^\top \lambda
\quad
\text{s.t.}\quad
\lambda \in \mathcal{K}_\mu,
\end{equation}
其中 $\mathcal{K}_\mu$ 表示由所有积分点局部库仑锥组成的直积锥，$H$ 与 $g$ 来自局部相对速度、局部线性化几何与近场接触响应。解得 $\lambda^\star$ 后，定义参考 wrench
\begin{equation}
 w_{\ell r}^\star
 =
 \sum_{k} G(x_k)\lambda_k^\star.
\end{equation}

这里的局部 dense 微求解并不替代全局 NSC，而只是为 reduced support set 提供当前时刻的高精参考目标。

\subsubsection{Reduced 力分配}

给定优化后的代表点集合 $Q_{\ell r}$ 与参考 wrench $w_{\ell r}^\star$，求解 reduced 点力：
\begin{equation}
\min_{\hat{\lambda}}
\frac{1}{2}\hat{\lambda}^\top W \hat{\lambda}
\end{equation}
\begin{equation}
\text{s.t.}\quad
\sum_{j=1}^{m} G(q_j)\hat{\lambda}_j = w_{\ell r}^\star,
\end{equation}
\begin{equation}
\hat{n}_j \ge 0,
\qquad
\sqrt{\hat{t}_{j1}^2+\hat{t}_{j2}^2}\le \mu_j \hat{n}_j.
\end{equation}

该问题同样是一个小型 SOCP。若可行，则说明当前 reduced support set 能对该 subpatch 的当前参考 wrench 实现\textbf{精确等效}；若不可行，则说明当前基不够强，必须执行以下操作之一：
\begin{enumerate}[label=(\arabic*)]
    \item 增加代表点数 $m$；
    \item 扩大代表点支撑域并重新优化；
    \item 将当前 subpatch 继续分裂。
\end{enumerate}

本文不允许在 reduced 力分配阶段以牺牲精度的方式“硬拟合”不可行解，而是始终以误差控制驱动模型升阶或细分。

\subsection{连续区域内部非穿透监测}

少量代表点接触方法的一个典型问题是：虽然代表点满足局部非穿透，但 patch 内部可能仍发生漏穿透。为此，本文在每个 subpatch 内布置一组不进入求解器的 sentinel 点
\begin{equation}
 \mathcal{S}_{\ell r} = \{s_i\}_{i=1}^{N_s},
\end{equation}
并在每一步预测后计算
\begin{equation}
 e_g
 =
 \max_{s_i\in\mathcal{S}_{\ell r}}
 \max(0,-g(s_i)).
\end{equation}

若 $e_g$ 超过阈值 $\varepsilon_g$，则说明 reduced support set 虽然在离散支撑点上满足约束，但未能充分覆盖连续接触区域的内部；此时优先对当前 subpatch 继续细分，并在必要时增加 support points。

\subsection{时间连续性与 Warm Start}

由于接触区域在相邻时间步之间通常具有时序连续性，本文为每个 subpatch 维护 persistent identifier，并在相邻步之间进行匹配：
\begin{equation}
\mathrm{match}(\Omega_{\ell r}^{t}, \Omega_{\ell' r'}^{t+\Delta t}).
\end{equation}

匹配指标包括质心距离、平均法向夹角、面积重叠率、局部骨架重叠与代表点集合的 Hausdorff 距离。对匹配成功的 subpatch，将上一时间步的 reduced impulses $\hat{\lambda}^{\mathrm{old}}$ 传输至新的 support set，传输过程同样采用保持 wrench 等效的优化：
\begin{equation}
\min_{\hat{\lambda}^{\mathrm{new}}}
\|\hat{\lambda}^{\mathrm{new}}\|_W^2
\quad
\text{s.t.}\quad
\sum_j G(q_j^{\mathrm{new}})\hat{\lambda}_j^{\mathrm{new}}
=
\sum_j G(q_j^{\mathrm{old}})\hat{\lambda}_j^{\mathrm{old}}.
\end{equation}

这意味着 warm start 的传递对象不是逐点冲量值本身，而是对应的等效 wrench。

\subsection{自适应误差控制}

综合前述定义，本文在每个 subpatch 上维护以下误差：
\begin{equation}
 e_{\mathrm{plane}},\quad
 e_{\mathcal{C}},\quad
 e_w=
 \frac{\|w^\star-\hat{w}\|}{\|w^\star\|+\varepsilon},\quad
 e_{\mathrm{CoP}},\quad
 e_g,\quad
 e_{\mathrm{topo}}.
\end{equation}
其中 $e_w$ 衡量参考 wrench 与 reduced wrench 的匹配误差，$e_{\mathrm{CoP}}$ 衡量 CoP 的漂移误差。自适应策略如下：
\begin{itemize}[leftmargin=2em]
    \item 若 $e_{\mathrm{plane}}$ 或 $e_{\mathrm{topo}}$ 主导，则细分 subpatch；
    \item 若 $e_{\mathcal{C}}$ 或 $e_{\mathrm{CoP}}$ 主导，则重新优化代表点位置；
    \item 若 reduced 力分配不可行，则增加代表点数；
    \item 若 $e_g$ 主导，则细分 subpatch 并增加 sentinel 点。
\end{itemize}

由此形成如下误差驱动循环：
\begin{enumerate}[label=\arabic*.]
    \item 构造 dense patch；
    \item patch / subpatch 分解；
    \item 代表点初始化与优化；
    \item 局部参考 wrench 预测；
    \item reduced SOCP 力分配；
    \item 误差评估；
    \item 若误差超阈值，则升阶或细分并重复。
\end{enumerate}

该过程体现了本文“\textbf{精度优先}”的设计原则：允许 reduced support set 的规模自适应增加，但不允许在误差失控的条件下维持低维表示。

\subsection{与 Chrono NSC 的集成}

Chrono 的 \verb|ChSystem| 提供了 \verb|SetContactContainer()| 接口，可替换底层 contact container；该 contact container 负责从碰撞检测系统收集接触信息并用于接触力生成。默认 \verb|AddContactCallback| 仅允许在接触被加入容器时修改复合材料参数，而不能替代几何级接触流形重构。因此，本文将所提方法集成为一个自定义 NSC contact container 或自定义接触报告模块，其工作流程为：Bullet 仅负责 broadphase 与碰撞对管理；对每个候选接触对，由本文方法构造 dense SDF 接触区域、执行多 patch / 多 subpatch 约化，并将所得 reduced support points 重新封装为普通 NSC 点接触，最终交由 Chrono NSC 的全局描述器和互补求解器处理。

设最终全部 subpatch 的代表点总数为
\begin{equation}
N_{\mathrm{red}} = \sum_{\ell=1}^{N_p}\sum_{r=1}^{m_\ell} m_{\ell r}.
\end{equation}

则系统级接触自由度规模由原始 dense 模型的 $O(N_q)$ 降至 $O(N_{\mathrm{red}})$，同时由于每个代表点仍以标准点接触形式存在，本文方法可以无缝兼容现有 NSC 的约束装配与互补求解过程。

\subsection{算法流程}

\noindent\textbf{算法 1：多 Patch / 多 Subpatch 保持 Wrench 等效的 SDF 接触约化}

\noindent 输入：候选接触对 $(A,B)$、SDF 几何、相对状态、误差阈值。\\
\noindent 输出：可注入 Chrono NSC 的 reduced point contacts。

\begin{enumerate}[label=\arabic*.]
    \item 由 SDF 构造高密积分点集合 $\mathcal{Q}$；
    \item 根据连通性与法向一致性得到 patch 集合 $\{\Omega_\ell\}$；
    \item 对每个 patch：
    \begin{enumerate}[label*=\arabic*.]
        \item 初始化 subpatch 队列；
        \item 对每个 subpatch：
        \begin{enumerate}[label*=\arabic*.]
            \item 计算局部平面化误差、二阶矩误差与拓扑误差；
            \item 若误差超阈值，则分裂 subpatch；
            \item 否则初始化代表点 $Q_{\ell r}$；
            \item 优化 $Q_{\ell r}$ 以减小 wrench 锥误差；
            \item 通过局部 dense 微求解得到参考 wrench $w_{\ell r}^\star$；
            \item 解 reduced SOCP，得到代表点力 $\hat{\lambda}$；
            \item 评估 $e_w,e_{\mathcal{C}},e_{\mathrm{CoP}},e_g$；
            \item 若任一误差超阈值，则升阶或细分并返回步骤 3.2.1。
        \end{enumerate}
    \end{enumerate}
    \item 将全部 reduced support points 转换为普通 NSC contact items；
    \item 调用 Chrono NSC 接触容器完成系统装配与互补求解。
\end{enumerate}

\subsection{方法讨论}

本文方法并不试图恢复真实的局部压力场，而是以 dense SDF 接触区域为参考，在\textbf{刚体动力学层面}保持其可行 wrench 锥与当前参考 wrench。相较于基于少量接触点缓存的流形方法，本文方法将接触约化从“缓存少量点”提升为“压缩一个连续接触区域在 wrench 空间中的可行集”；相较于直接采用 patch-level 六维接触元的方法，本文方法又保留了与 Chrono NSC 标准点接触求解链路的兼容性。因此，它特别适用于以下场景：需要高精度 body-level 合力、合力矩与作用线等效，但又希望继续复用现有 NSC 互补求解器的刚体接触动力学问题。

\begin{remark}
本文方法的核心边界在于：它追求的是刚体层面的高精度 wrench 等效，而非局部压力场、接触应力峰值或真实材料变形场的逐点重建。因此，当研究目标转向材料级接触应力与弹性变形主导的接触斑演化时，更适合采用分布式压力场或柔顺接触建模框架。
\end{remark}

\section{理论分析与性质讨论}
\label{sec:analysis}

\subsection{单点纯力等效的边界}

设某连续接触区域关于参考点 $O$ 的结果 wrench 为 $(F,M_O)$。若试图仅用一个作用于点 $q$ 的单力 $F$ 来等效，则必有
\begin{equation}
M_O = (q-O)\times F.
\end{equation}
因此必然满足
\begin{equation}
M_O \cdot F = 0.
\end{equation}

这说明：当连续接触区域存在沿合力方向的自由矩分量时，单点纯力表示不再精确可行，必须引入多个接触点或 patch-level 接触 wrench。该结果也从理论上解释了为何仅靠单个代表点无法统一处理一般面接触与长条接触。

\subsection{局部共面条件下的三点完备性}

局部共面 subpatch 的三点完备性已在第~\ref{sec:method} 节给出。进一步地，对任意目标局部 wrench $w=(T_x,T_y,N,M_x,M_y,M_z)^\top$，若三点非共线，则在忽略库仑锥约束时，总存在一点力组合使三点系统精确表示该 wrench。换言之，三点支持集已经构成了局部六维 wrench 的最小生成基。

这一性质为本文的多级约化策略提供了基础：默认情况下，局部近共面 subpatch 以三点表示即可；仅当法向可行性、摩擦可行性、CoP 支撑范围或内部 gap 误差不满足时，才需要增加点数或细分 subpatch。

\subsection{局部平面化误差的定性上界}

设某 subpatch 的最大主曲率为 $\kappa_{\max}$，其局部直径为 $D$。若以局部切平面对其进行几何近似，则高度误差具有典型的二阶尺度，可记为
\begin{equation}
\Delta h = O(\kappa_{\max} D^2).
\end{equation}

对应到关于 patch 质心的力矩误差，其量纲上可写为
\begin{equation}
\|\Delta M\| \le C\,\kappa_{\max}D^2\|F\|,
\end{equation}
其中 $C$ 为与局部参数化和载荷分布有关的常数。该表达式并非严格封闭形式上界，而是说明：当曲率增大或 subpatch 尺度增大时，局部共面假设将迅速变差，必须通过细分 subpatch 控制误差。

\subsection{摩擦可行性与代表点支撑半径}

对给定总法向载荷 $N$ 和切向扭矩需求 $M_z$，若 reduced support set 的几何支撑半径过小，则即使总切向合力大小不大，仍可能因杠杆臂不足而无法满足
\begin{equation}
\sqrt{\hat{t}_{j1}^2+\hat{t}_{j2}^2}\le \mu_j\hat{n}_j.
\end{equation}

因此，代表点的位置不仅影响法向 CoP 的可行性，也直接影响切向扭矩的可实现性。这一现象说明，仅通过“选择最深的几点”并不能构造好的支持基；需要同时考虑几何尺度、一二阶矩与 wrench cone gap。

\subsection{复杂度讨论}

设某接触对的 dense 积分点数为 $N_q$，最终 reduced 代表点总数为 $N_{\mathrm{red}}$。则本文方法的额外代价主要来自：
\begin{enumerate}[label=(\arabic*)]
    \item dense 接触区域采样与 patch 图构建；
    \item patch / subpatch 自适应细分；
    \item 每个 subpatch 的小规模代表点优化；
    \item 每个 subpatch 的局部 dense 微求解与 reduced SOCP 力分配。
\end{enumerate}

相较于直接将全部 dense 积分点送入全局 NSC 求解器，本文方法将全局约束维度控制在 $O(N_{\mathrm{red}})$，并把精度主要花在局部 subpatch 的小规模优化上。这一设计适合“局部高精、全局保持可解”的仿真场景。

\section{实验设计}
\label{sec:experiments}

本节给出面向后续投稿的实验设计框架。由于当前版本重点在方法论与理论层面，因此实验部分先给出完整设计、评价指标和对比基线，具体数值结果可在实现完成后补入。

\subsection{实验目标}

实验的总体目标是验证：本文方法相较于传统点接触流形维护方式，是否能够在保持 Chrono NSC 全局求解链路不变的前提下，更准确地重建连续接触区域在刚体动力学层面的等效 wrench，并在不同接触类型下保持稳定和一致的高精度表现。

具体而言，实验应回答以下问题：
\begin{enumerate}[label=(\arabic*)]
    \item 在点接触、连续区域接触和长条形接触场景中，本文方法是否能显著降低合力和合力矩误差？
    \item 在接触作用线、CoP 漂移和内部 gap 误差上，本文方法是否优于传统流形维护方法？
    \item 多 patch / 多 subpatch 的自适应 refinement 是否能根据接触类型自动调节模型复杂度？
    \item 在不同时间步长、摩擦系数和几何分辨率下，本文方法是否具有较好的鲁棒性？
\end{enumerate}

\subsection{实验平台与实现设置}

建议的实验平台如下：
\begin{itemize}[leftmargin=2em]
    \item 动力学后端：Chrono NSC；
    \item 广相与候选对管理：Bullet broadphase；
    \item 几何表征：SDF 或表面导出 SDF；
    \item dense 参考模型：高密度积分点 + 局部微求解；
    \item reduced 模型：本文提出的多 patch / 多 subpatch wrench-preserving reduction；
    \item 对比模型：默认 NSC 点接触流形、固定点数 manifold reduction、无 subpatch 的单 patch 代表点模型。
\end{itemize}

为保证可重复性，应统一以下参数：时间步长、求解器容差、摩擦系数、SDF 分辨率、接触采样密度、误差阈值和最大代表点数等。

\subsection{测试场景设计}

\subsubsection{场景 A：点接触基准测试}

该场景用于验证本文方法在典型点接触问题中不会引入不必要误差。建议采用如下模型：
\begin{itemize}[leftmargin=2em]
    \item 小球--小球对心碰撞；
    \item 小球--平面弹跳；
    \item 低速滚动进入单点碰撞。
\end{itemize}

该类场景下，dense 接触区域本身很小，理论上本文方法应当退化为极少量 patch 与极少量代表点，并与标准点接触结果基本一致。重点评价该方法是否保持原有点接触问题的准确性与数值稳定性。

\subsubsection{场景 B：连续区域接触测试}

该场景用于检验本文方法在面接触或连续区域接触中的 wrench 等效能力。建议采用如下模型：
\begin{itemize}[leftmargin=2em]
    \item 凸轮--从动件接触；
    \item 平板--平板局部压靠；
    \item 曲面块体沿导轨表面滚滑接触。
\end{itemize}

该类场景中，重点观察本文方法是否能正确保持合力、合力矩和接触作用线，尤其是法向 CoP 的稳定性与切向摩擦矩的再现能力。

\subsubsection{场景 C：长条形接触测试}

该场景用于检验多 subpatch 策略的必要性和有效性。建议采用如下模型：
\begin{itemize}[leftmargin=2em]
    \item 齿轮齿面啮合局部接触；
    \item 插销--轴承或轴--孔接触；
    \item 细长圆柱与曲面导轨接触。
\end{itemize}

这类场景的典型特点是接触沿一个主方向强烈延展，且局部法向沿接触线发生变化。对其进行单 patch 约化通常会造成明显力矩误差，因此该类实验特别适合验证本文的多 subpatch 自适应分解策略。

\subsubsection{场景 D：多 patch / 不连通接触测试}

建议构造两个或多个不连通 patch 同时接触的情形，例如：
\begin{itemize}[leftmargin=2em]
    \item 非凸零件与平板的多点支撑；
    \item 带凹槽零件的双边接触；
    \item 多凸起接触足与地面的同时接触。
\end{itemize}

该场景用于验证本文方法是否能在拓扑上正确区分多个 patch，并避免不同接触区域被错误合并。

\subsection{评价指标}

建议采用如下指标：

\paragraph{(1) Wrench 误差}
定义 dense 参考模型给出的 wrench 为 $w_{\mathrm{dense}}$，reduced 模型给出的 wrench 为 $w_{\mathrm{red}}$，则
\begin{equation}
\varepsilon_w = \frac{\|w_{\mathrm{dense}}-w_{\mathrm{red}}\|}{\|w_{\mathrm{dense}}\|+\epsilon}.
\end{equation}

\paragraph{(2) 合力误差与合力矩误差}
可分别定义
\begin{equation}
\varepsilon_F = \frac{\|F_{\mathrm{dense}}-F_{\mathrm{red}}\|}{\|F_{\mathrm{dense}}\|+\epsilon},
\qquad
\varepsilon_M = \frac{\|M_{\mathrm{dense}}-M_{\mathrm{red}}\|}{\|M_{\mathrm{dense}}\|+\epsilon}.
\end{equation}

\paragraph{(3) CoP 误差}
若法向合力非零，则定义
\begin{equation}
\varepsilon_{\mathrm{CoP}} = \|c_{\mathrm{cop}}^{\mathrm{dense}}-c_{\mathrm{cop}}^{\mathrm{red}}\|.
\end{equation}

\paragraph{(4) 内部 gap 误差}
\begin{equation}
\varepsilon_g = \max_{s_i\in\mathcal{S}} \max(0,-g(s_i)).
\end{equation}
该指标用于反映 reduced 点接触是否在 patch 内部遗漏了非穿透约束。

\paragraph{(5) 拓扑保持指标}
对多个 patch 场景，统计 reduced 模型恢复的 patch 个数与 dense 参考模型的 patch 个数是否一致，并记录错误合并与错误分裂次数。

\paragraph{(6) 求解规模与运行时间}
记录最终代表点总数 $N_{\mathrm{red}}$、平均 subpatch 数、每步局部微求解时间、全局求解时间与总仿真时间，用于评估精度与计算代价之间的关系。

\subsection{对比基线}

建议至少设置以下基线：
\begin{enumerate}[label=(\arabic*)]
    \item \textbf{Baseline-1：}Chrono NSC 默认点接触管线；
    \item \textbf{Baseline-2：}固定点数的接触流形约化，如直接保留少量接触点；
    \item \textbf{Baseline-3：}单 patch 的代表点模型，不进行多 subpatch 细分；
    \item \textbf{Baseline-4：}本文方法的消融版本，例如去掉局部参考 wrench 预测、去掉 sentinel gap 监测或去掉 CoP 误差约束。
\end{enumerate}

通过这些基线，可分别回答“多 patch 是否必要”“多 subpatch 是否必要”“wrench-preserving 是否必要”“局部 dense 微求解是否必要”等关键问题。

\subsection{消融实验设计}

建议进行如下消融：
\begin{enumerate}[label=(\alph*)]
    \item 去掉 $e_{\mathcal{C}}$，仅保留几何矩匹配，观察 wrench 误差变化；
    \item 去掉 CoP 约束或 CoP 误差驱动，观察法向力矩与作用线漂移；
    \item 去掉内部 sentinel 监测，观察 patch 内部是否出现漏穿透；
    \item 将每个 subpatch 的代表点数固定为 3，不允许升阶，观察复杂接触下的误差上升；
    \item 用单 patch 替代多 subpatch，观察长条形接触下的误差变化。
\end{enumerate}

\subsection{结果展示建议}

正式论文中建议采用以下图表：
\begin{itemize}[leftmargin=2em]
    \item 不同场景下 dense patch 与 reduced support set 的可视化图；
    \item 不同时间步的 patch / subpatch 分解与跟踪结果；
    \item 合力误差、合力矩误差、CoP 误差与内部 gap 误差的折线图或箱线图；
    \item 不同误差阈值下 $N_{\mathrm{red}}$ 与精度的关系曲线；
    \item 消融实验表格与复杂度统计表。
\end{itemize}

\subsection{预期结果与讨论重点}

本文预期结果包括：
\begin{enumerate}[label=(\arabic*)]
    \item 在点接触场景中，本文方法不显著劣于默认点接触模型；
    \item 在连续区域接触与长条形接触中，本文方法显著降低合力矩误差与 CoP 漂移；
    \item 多 subpatch 机制对长条形接触和多 patch 接触具有决定性作用；
    \item 在相同全局求解框架下，本文方法能以适中的代表点数显著提升刚体层面的接触等效精度。
\end{enumerate}

讨论时应特别强调：本文方法提升的是 \textbf{body-level wrench 等效精度}，而非逐点压力场精度。因此，实验结果应围绕动力学层面的量展开，而不应把其与完整连续接触压力场重建混为一谈。

\section{结论与未来工作}
\label{sec:conclusion}

本文围绕 Chrono NSC 在连续区域接触高精度建模中的不足，提出了一种面向 SDF 的多 patch / 多 subpatch 保持 wrench 等效的接触约化方法。该方法以 dense SDF 接触区域为参考，通过 patch / subpatch 自适应分解、代表点优化、局部参考 wrench 预测、reduced SOCP 力分配以及内部 gap 监测，构建出一种既面向精度、又兼容现有 NSC 点接触求解链路的低维接触表示。本文的核心观点是：接触约化不应被视为简单的几何删点问题，而应被视为从 dense admissible wrench cone 到 reduced admissible wrench cone 的误差控制投影问题。

从理论层面看，本文给出了局部共面 subpatch 的三点 wrench 完备性分析，并建立了法向可行性与 CoP 支撑域之间的几何联系；从工程实现层面看，本文方法提供了一条不重写 Chrono 全局互补求解框架、但显著改善连续接触区域刚体层面等效精度的可落地路径。

未来工作可沿以下方向展开：
\begin{enumerate}[label=(\arabic*)]
    \item 将当前方法从刚体层面的 wrench 等效扩展到更强的 patch-level 接触元表示；
    \item 结合更高效的 SDF 数据结构与局部微求解器，提高大规模场景中的实时性；
    \item 在柔顺接触、弹性体接触与可微仿真框架中研究类似的 patch-level reduction 机制；
    \item 将该方法系统实现于 Chrono 后端，并与默认 NSC 管线进行大规模定量对比。
\end{enumerate}

\appendix
\section{符号表（占位）}

\begin{longtable}{p{0.2\textwidth}p{0.7\textwidth}}
\toprule
符号 & 含义 \\
\midrule
$\phi_A,\phi_B$ & 刚体 $A,B$ 的符号距离场 \\
$g(x)$ & 局部 gap 场 \\
$\Omega$ & 连续接触区域 \\
$\Omega_\ell$ & 第 $\ell$ 个 patch \\
$\Omega_{\ell r}$ & 第 $\ell$ 个 patch 的第 $r$ 个 subpatch \\
$\mathcal{Q}$ & dense 积分点集合 \\
$Q_{\ell r}$ & reduced support set \\
$w$ & 局部六维 wrench \\
$\mathcal{C}_{\mathrm{dense}}$ & dense 可行 wrench 锥 \\
$\mathcal{C}_{\mathrm{red}}$ & reduced 可行 wrench 锥 \\
$e_{\mathrm{plane}}$ & 局部平面化误差 \\
$e_{\mathcal{C}}$ & wrench 锥误差 \\
$e_w$ & 参考 wrench 匹配误差 \\
$e_{\mathrm{CoP}}$ & 中心压力点误差 \\
$e_g$ & 内部 gap 误差 \\
\bottomrule
\end{longtable}

\section*{参考文献（占位）}
\begin{thebibliography}{99}

\bibitem{chrono_system}
Radu Serban, Alessandro Tasora, Dan Negrut, et al.
\newblock Project Chrono: An open-source multi-physics dynamics engine.
\newblock [占位：请替换为正式 Chrono 论文或官方文档条目].

\bibitem{chrono_nsc}
Project Chrono Documentation.
\newblock Non-smooth contact and system architecture.
\newblock [占位：请替换为 Chrono NSC 与 contact container 的正式文献或官方 API 文档].

\bibitem{bullet_manifold}
Erwin Coumans.
\newblock Bullet Physics Library and persistent contact manifold.
\newblock [占位：请替换为 Bullet 官方论文、文档或源码引用条目].

\bibitem{drake_hydro}
Russ Tedrake, et al.
\newblock Hydroelastic contact model in Drake.
\newblock [占位：请替换为 Drake hydroelastic 相关论文或官方文档].

\bibitem{ecp_contact}
Mohamed A. Posa, or related authors.
\newblock Equivalent contact point and contact wrench representations for line/surface contact.
\newblock [占位：请替换为 ECP 相关正式论文条目].

\bibitem{sdf_contact}
[作者占位].
\newblock SDF-based contact modeling for rigid or compliant simulation.
\newblock [占位：补充与你方案最贴近的 SDF 接触论文].

\bibitem{contact_manifold_reduction}
[作者占位].
\newblock Contact manifold simplification and reduction in multibody dynamics or graphics simulation.
\newblock [占位：补充接触流形约化相关论文].

\bibitem{frictional_contact}
[作者占位].
\newblock Frictional contact complementarity, variational inequality, and numerical solution methods.
\newblock [占位：补充 NSC / 摩擦接触求解理论的基础参考文献].

\end{thebibliography}

\end{document}
